\section{Introduction}
\label{sec:introduction}

Arithmetic symbolic dynamics faces a recurrent design tension.  A transition
grammar can be sufficiently rigid to encode a desired arithmetic list, but
then the recurrence may merely replay a verifier.  Conversely, an intrinsic
recurrent grammar may possess a robust determinant while producing far too
many primitive cycles.  The successor--divisor shift sits on the second side
of this divide.  Its vertex \(n\) represents the full shift on \(n\) symbols,
and its edge rule

\[
 n\longrightarrow d
 \quad\Longleftrightarrow\quad d\ge2,\quad d\mid n+1
\]

uses only alphabet successor and tensor factorization.  The resulting
countable Markov graph is recurrent and its endpoint-weighted adjacency has a
sharp trace-class half-plane.  It also contains the simple cycle

\[
 C_k=(k,k+1,\ldots,2k-1)
\]

for every \(k\ge2\).  The analytic determinant is therefore real structure,
but its primitive species is not selective.

This paper asks whether information already present on an allowed edge can
resolve that flood.  The edge equation \(n+1=dq\) exposes the unique
cofactor

\[
 q(n,d)=\frac{n+1}{d}.
\]

No rational-prime table or target spectral data are needed to define it.
Multiplying these labels around a closed path produces a canonical positive
monoid cocycle.  Passing to the group completion gives regular extensions,
unitary character fibers, and exact Fourier coefficients of a connected
Fredholm ledger.  These are standard kinds of constructions in dynamical
zeta and graph-cover theory
\citep{BowenLanford1970,AdachiSunada1987,GrossTucker1977,StarkTerras1996};
the question is what this particular source-derived cocycle actually knows.

The answer has a strong positive part.  The cycle product telescopes to

\[
 Q(\gamma)=\prod_{n\in\gamma}\left(1+\frac1n\right)>1.
\]

The class \(Q=2\) consists exactly of the canonical cycles \(C_k\), one
primitive class at every length.  More generally, an atomic holonomy \(p\)
has exactly the representatives

\[
 C_{k,p}=(k,k+1,\ldots,pk-1),\qquad k\ge2.
\]

Temporal repetitions cannot enter an atomic class.  Thus the character
family resolves a nontrivial arithmetic grading exactly, rather than merely
decorating an uncontrolled orbit sum.

The analytic part is equally sharp.  On \(\ell^2\{2,3,\ldots\}\), define

\[
 L_{s,u}e_n=
 \sum_{d\mid n+1,\ d\ge2}
 (nd)^{-s}q(n,d)^{-u}e_d.
\]

We prove

\[
 \boxed{
 L_{s,u}\in\Sone
 \iff
 \Re s>\frac12,
 \quad \Re(s+u)>\frac12.}
\]

The two conditions are independent.  Output-row tails detect the combined
endpoint--cofactor decay.  A contractive Fourier projection detects the
\(q=1\) successor spine, on which the cofactor twist disappears completely.
This gives a full phase diagram, not a finite-cutoff heuristic.

The same theorem also identifies the obstruction.  At \(s=0\), the pure
cofactor roof would assign the attractive cycle weight \(Q^{-u}\), but its
successor component is the unweighted unilateral shift.  The matrix is never
trace class and is noncompact whenever bounded.  Restoring endpoint decay
produces an honest Fredholm determinant but assigns the \(Q=2\) cycle the
factorial weight

\[
 \left(\frac{(2k-1)!}{(k-1)!}\right)^{-2s}.
\]

Unitary characters supply phases only: all \(C_k\) carry the same phase
\(\chi(2)\).  We call this the \emph{Fredholm trilemma}.

The regular group lift sharpens the distinction.  Relative to its natural
semifinite group trace it is \(L^1\) for \(\Re s>1/2\), yet it is not an
ordinary compact operator because of infinite deck multiplicity.  Moreover,
the neutral group trace sees no periodic path and its normalized local
determinant is exactly one.  Semifinite integrability, ordinary Fredholm
theory, and neutral recurrence are therefore three different statements.

Our contributions are:

\begin{enumerate}
\item an exact positive-holonomy and gauge theorem for the source cofactor;
\item complete \(Q=2\) and atomic-holonomy classifications, with exact
connected character coefficients;
\item the sharp two-parameter trace-class theorem and the separate regular
lift analysis;
\item a source-level no-go showing that abelian product holonomy is blind to
the index \(k\), stable under positive-inventory controls.
\end{enumerate}

These results are constructive but not a Riemann representation.  For every
atom \(p\), the graph produces an infinite grid \((p,k)\), not one primitive
orbit per rational prime, and its honest endpoint roofs are factorial.  The
strict evaluation is therefore

\[
\begin{aligned}
(&\mathrm{A0\_STRUCTURAL\_ARITHMETIC\_RELATION},\\
 &\mathrm{A1\_WEAK},\\
 &\mathrm{A2\_ANALYTIC\_DETERMINANT},\\
 &\mathrm{A3\_FAIL},\\
 &\mathrm{A4\_FAIL}).
\end{aligned}
\]

with \(\mathrm{ROUTE\_A\_REJECTED}\).  No target-zero data enter any theorem,
and Route B remains locked.

\paragraph{Organization.}
\Cref{sec:classical} fixes the literature boundary.
\Cref{sec:source} defines the symbolic object and determinant conventions.
\Cref{sec:holonomy,sec:atomic} prove the algebraic classification.
\Cref{sec:ledger} resolves connected coefficients.
\Cref{sec:traceclass,sec:lift} prove the scalar and lifted operator theorems.
\Cref{sec:trilemma,sec:controls} give the no-go and controls, and
\Cref{sec:route} records the route decision.
